% Figure: laminar (parabolic) versus turbulent (1/7-power, blunt) pipe
% velocity profiles at the same mean velocity. The turbulent profile is
% flatter across the core and steeper at the wall, which is why it carries
% a larger wall shear stress and a larger friction factor.
\begin{figure}[htbp]
\centering
\begin{tikzpicture}
\begin{axis}[
  width=11cm, height=7cm,
  xlabel={$u_z / u_{\max}$}, ylabel={$r/R$},
  xmin=0, xmax=1.1, ymin=-1.05, ymax=1.05,
  axis lines=left,
  legend pos=outer north east,
  legend cell align=left,
  grid=major, grid style={enggray!20},
  tick label style={font=\scriptsize},
  label style={font=\small},
  legend style={font=\scriptsize},
]
% Laminar parabola u = 1-(r/R)^2.
\addplot[engblue, very thick, domain=-1:1, samples=160] ({1 - x*x}, {x});
\addlegendentry{laminar (parabolic)}
% Turbulent 1/7-power: u = (1 - |r/R|)^{1/7}. Clamp argument away from
% negative by taking abs; sample upper and lower halves separately.
\addplot[engred, very thick, domain=0:1, samples=120]
  ({(1 - x)^(1/7)}, {x});
\addplot[engred, very thick, domain=0:1, samples=120, forget plot]
  ({(1 - x)^(1/7)}, {-x});
\addlegendentry{turbulent ($1/7$-power)}
% pipe walls
\addplot[engblack, very thick] coordinates {(0,1) (1.1,1)};
\addplot[engblack, very thick] coordinates {(0,-1) (1.1,-1)};
\end{axis}
\end{tikzpicture}
\caption[Laminar versus turbulent pipe profiles]{Laminar and turbulent
pipe velocity profiles, each normalised to its own centreline maximum.
The turbulent profile (the empirical $1/7$-power law) is blunt across the
core and falls steeply through a thin near-wall region, so it carries a
much larger wall velocity gradient and a larger friction factor than the
laminar parabola at the same Reynolds number.}
\label{fig:vol03ch12:laminar-turbulent}
\end{figure}
